%% ============================================================
%% preamble/styling.tex — section format, running heads, spacing, draft
%%
%% Purpose      : how the paper LOOKS — titlesec section formats,
%%                fancyhdr running heads, line spacing (setspace), the
%%                theorem environments, widow/orphan policy, and the
%%                \DraftMode build-mode handling.
%% Dependencies : titlesec + fancyhdr + setspace + amsthm (packages.tex),
%%                colors.tex (pt-* tokens), generated/metadata.tex.
%% Provenance   : titlesec section rules in the primary color, from the
%%                datacenter/moratorium SSRN-cluster styling.tex; the
%%                DRAFT banner discipline from book-template styling.tex.
%% ============================================================

\makeatletter

% ----------------------------------------------------------------------
% Line spacing (paper.yaml typography.linespacing). Law reviews double.
% ----------------------------------------------------------------------
\ifdefstring{\PaperLineSpacing}{double}{\doublespacing}{%
  \ifdefstring{\PaperLineSpacing}{onehalf}{\onehalfspacing}{\singlespacing}}

% ----------------------------------------------------------------------
% ISO date for the draft banner.
% ----------------------------------------------------------------------
\newcommand{\pt@isodate}{\the\year-\two@digits\month-\two@digits\day}

% ----------------------------------------------------------------------
% Build modes. \DraftMode: line numbers + a "DRAFT" banner (never a
% diagonal watermark — it fights running heads). The default build
% suppresses it. \AnonMode is read in frontmatter/titleblock.
% ----------------------------------------------------------------------
\ifdefined\DraftMode
  \usepackage[switch,columnwise]{lineno}
  \newcommand{\pt@draftbanner}{%
    {\normalfont\sffamily\bfseries\footnotesize\color{pt-draft}%
      DRAFT\,\textperiodcentered\,\pt@isodate\,\textperiodcentered\,Comments welcome}}
\else
  \newcommand{\pt@draftbanner}{}
\fi

% ======================================================================
% Section headings — titlesec, sans, primary color, with a thin rule
% under \section (the SSRN-cluster look).
% ======================================================================
\titleformat{\section}
  {\normalfont\Large\sffamily\bfseries\color{pt-section}}
  {\thesection}{0.6em}{}[{\color{pt-section-rule}\titlerule[0.6pt]}]
\titleformat{\subsection}
  {\normalfont\large\sffamily\bfseries\color{pt-section}}
  {\thesubsection}{0.6em}{}
\titleformat{\subsubsection}
  {\normalfont\normalsize\sffamily\bfseries\color{pt-section}}
  {\thesubsubsection}{0.6em}{}
\titlespacing*{\section}{0pt}{1.4\baselineskip}{0.6\baselineskip}
\titlespacing*{\subsection}{0pt}{1.1\baselineskip}{0.4\baselineskip}

% ======================================================================
% Running heads — short title (recto) + page number (outer). Draft mode
% replaces the centered head with the DRAFT banner.
% ======================================================================
\pagestyle{fancy}
\fancyhf{}
\renewcommand{\headrulewidth}{0pt}
\ifdefined\DraftMode
  \fancyhead[C]{\pt@draftbanner}
  \fancyhead[R]{\color{pt-folio}\thepage}
\else
  \fancyhead[L]{\color{pt-running-head}\footnotesize\PaperShortTitle}
  \fancyhead[R]{\color{pt-folio}\thepage}
\fi
% First page of each section (\section* title pages fall back to plain);
% keep the folio in the footer there.
\fancypagestyle{plain}{%
  \fancyhf{}\renewcommand{\headrulewidth}{0pt}%
  \fancyhead[C]{\pt@draftbanner}%
  \fancyfoot[C]{\color{pt-folio}\thepage}}

% ======================================================================
% Theorem environments (amsthm) — shared numbering by section.
% ======================================================================
\theoremstyle{plain}
\newtheorem{theorem}{Theorem}[section]
\newtheorem{lemma}[theorem]{Lemma}
\newtheorem{proposition}[theorem]{Proposition}
\newtheorem{corollary}[theorem]{Corollary}
\theoremstyle{definition}
\newtheorem{definition}{Definition}[section]
\newtheorem{assumption}{Assumption}
\theoremstyle{remark}
\newtheorem{remark}{Remark}

% ======================================================================
% Float + caption polish; page-breaking policy.
% ======================================================================
\captionsetup{labelfont={bf,color=primary}, textfont={color=text-primary}}
\captionsetup[sub]{labelfont={bf}, textfont={}}

\widowpenalty=10000
\clubpenalty=10000
\displaywidowpenalty=10000
\interfootnotelinepenalty=10000
\tolerance=1500
\emergencystretch=1.5em

\makeatother
